\documentclass[11pt]{amsart}

\usepackage[T1]{fontenc}
\usepackage{lmodern}
\usepackage{microtype}
\usepackage{amsmath,amssymb}
\usepackage[hidelinks]{hyperref}

\newtheorem{theorem}{Theorem}
\newcommand{\F}{\mathbb{F}}

\title[A $K_8$ counterexample]
      {A $K_8$ counterexample to Conjecture~3.9 of Bar\'at and Nagy}
\date{}

\subjclass[2020]{05C15, 05C70}
\keywords{rainbow subgraph, proper edge-coloring, 2-factor,
one-factorization}

\begin{document}

\begin{abstract}
Bar\'at and Nagy conjectured that every proper edge-coloring of $K_{2n}$
by at least $2n-1$ colors contains a rainbow $2$-regular subgraph on
$2n-1$ or $2n$ vertices.  We show that the conjecture fails for $n=4$.  The
coloring is the one-factorization represented by matrix $B$ in work of
Akbari, Etesami, Mahini, and Mahmoody, who reported that it has no
rainbow $7$-cycle.  We identify this earlier coloring as a counterexample
and give a self-contained algebraic proof that it has no rainbow
$2$-regular subgraph on seven vertices.
\end{abstract}

\maketitle

A subgraph is \emph{rainbow} if its edges have pairwise distinct colors.
We use ``$2$-factor on $m$ vertices'' for a $2$-regular subgraph with
$m$ vertices, as in \cite{BN}.  Conjecture~3.9 of Bar\'at and Nagy states
that every proper edge-coloring of $K_{2n}$ by at least $2n-1$ colors
contains a rainbow $2$-factor on $2n-1$ or $2n$ vertices.\footnote{The
statement is Conjecture~3.9 of the published version \cite{BN}; it appears
as Conjecture~3.8 of the preprint \cite{BNpre}.}  Since
$\chi'(K_{2n})=2n-1$, a proper $(2n-1)$-edge-coloring is a
one-factorization, so a one-factorization counterexample refutes the
``at least'' form as well.  We disprove the conjecture at $n=4$.

\begin{theorem}\label{thm:main}
There is a proper $7$-edge-coloring of $K_8$ with no rainbow $2$-regular
subgraph on seven or eight vertices.  Consequently, Conjecture~3.9 of
Bar\'at and Nagy is false.
\end{theorem}

\begin{proof}
Let
\[
  \Gamma=\F_2^2=\{0,a,b,a+b\},
\]
and let $\rho\colon\Gamma\to\Gamma$ be the linear automorphism defined by
\[
  \rho(a)=a+b,
  \qquad
  \rho(b)=a.
\]
Then $\rho(a+b)=b$, so $\rho$ cyclically permutes the three nonzero
elements of $\Gamma$ and has no nonzero fixed point.

Take $V(K_8)=\Gamma\times\{0,1\}$.  For $d\in\Gamma\setminus\{0\}$ assign
the internal edge $\{(x,i),(x+d,i)\}$ the color $I_d$, and assign the
cross edge $\{(x,0),(y,1)\}$ the color
\[
  C_{\rho(x)+y}.
\]
Each $I_d$ is a perfect matching.  Each $C_s$ is a perfect matching
because $x\mapsto s+\rho(x)$ is a bijection of $\Gamma$.  Thus the three
internal colors and four cross colors form a one-factorization of $K_8$.

Suppose that $F$ is a rainbow $2$-regular subgraph on seven vertices.
It has seven edges and hence uses every color exactly once.  In
particular, it has four cross edges and three internal edges.  Put
\[
  r_i=\lvert V(F)\cap(\Gamma\times\{i\})\rvert,
  \qquad
  m_i=\lvert E(F[\Gamma\times\{i\}])\rvert.
\]
Degree sums in the two parts give
\[
  2r_i=2m_i+4.
\]
Since $r_0+r_1=7$, one part contains three vertices and one internal edge,
while the other contains four vertices and two internal edges.  Let
$I_d$ be the color of the single internal edge and $I_e,I_f$ the colors
of the other two.  The three internal colors are all used, so
\[
  e+f=d.
\]

Write the cross edges, with incidence labels repeated when necessary, as
\[
  \{(x_j,0),(y_j,1)\},
  \qquad 1\le j\le4.
\]
In each part, the sum of the labels over all edge incidences is zero,
since every vertex of $F$ has degree two and $\Gamma$ has characteristic
two.  An internal edge of color $I_t$ has endpoint sum $t$.  The
internal-edge sum is therefore $d$ in both parts: it is $d$ on one side
and $e+f=d$ on the other.  Hence
\[
  \sum_{j=1}^4x_j=d,
  \qquad
  \sum_{j=1}^4y_j=d.
\]
All four cross colors occur once, and the sum of the elements of
$\Gamma$ is zero.  Therefore
\[
  0
  =\sum_{s\in\Gamma}s
  =\sum_{j=1}^4\bigl(\rho(x_j)+y_j\bigr)
  =\rho(d)+d,
\]
contradicting the absence of a nonzero fixed point of $\rho$.  Thus no
rainbow $2$-regular subgraph on seven vertices exists.

A $2$-regular subgraph on all eight vertices has eight edges and cannot
be rainbow in a coloring with only seven colors.
\end{proof}

\paragraph{Prior art and scope.}
Under the vertex order
\[
 (0,0),(a,0),(b,0),(a+b,0),(0,1),(a,1),(b,1),(a+b,1)
\]
and the color order
\[
 I_a,I_b,I_{a+b},C_0,C_a,C_b,C_{a+b},
\]
this is exactly matrix $B$ in \cite{AEMM}.  Akbari, Etesami, Mahini, and
Mahmoody reported, based on a computer search, that the coloring has no
rainbow $7$-cycle.  The one-factorization and that property are prior.
The proof above excludes both possible $2$-regular graphs on seven
vertices, $C_7$ and $C_3\cup C_4$, and records the resulting
counterexample to Conjecture~3.9.  The analogous statement restricted to
$n\ge5$ is not addressed here.

\begin{thebibliography}{9}

\bibitem{AEMM}
S.~Akbari, O.~Etesami, H.~Mahini, and M.~Mahmoody,
\emph{On rainbow cycles in edge colored complete graphs},
Australas. J. Combin. \textbf{37} (2007), 33--42.

\bibitem{BN}
J.~Bar\'at and Z.~L.~Nagy,
\emph{Transversals in generalized Latin squares},
Ars Math. Contemp. \textbf{16} (2019), no.~1, 39--47,
\url{https://doi.org/10.26493/1855-3974.1316.2d2}.

\bibitem{BNpre}
J.~Bar\'at and Z.~L.~Nagy,
\emph{Transversals in generalized Latin squares},
preprint, arXiv:1701.08220v1 (2017),
\url{https://arxiv.org/abs/1701.08220v1}.

\end{thebibliography}

\end{document}